% Example "kepler": the problem, the orbits against the exact ellipse, and the
% convergence and invariant tests. Every number comes from data/
% (parameters.tex, results.tex, ellipse.dat, orbit_<method>.dat,
% history_<method>.dat, error_<method>.dat, invariants_<method>.dat), written
% by examples/kepler/main.f90: run `fpm run --example kepler` from the project
% root, then build this file in examples/kepler/.
\documentclass[11pt,a4paper]{article}
\usepackage[margin=2.2cm]{geometry}
\usepackage{amsmath,booktabs}
\usepackage{pgfplots}
\usepgfplotslibrary{groupplots}
\pgfplotsset{compat=1.18}
\pagestyle{empty}

\begin{document}
\section*{Example: the Kepler problem}

\subsection*{1. Problem}
A body around a fixed mass, in units where $GM = 1$:
\begin{equation*}
  \ddot q = -\frac{q}{|q|^3}, \qquad
  q(0) = (1 - e,\, 0), \qquad
  \dot q(0) = \Bigl(0,\, \sqrt{\tfrac{1 + e}{1 - e}}\Bigr),
\end{equation*}
an ellipse of eccentricity $e$, semi-major axis 1 and period $2\pi$, started at
the pericentre, where it is fastest. The exact solution follows from Kepler's
equation for the eccentric anomaly $E(t)$:
\begin{equation*}
  E - e \sin E = t, \qquad
  q = \bigl(\cos E - e,\; \sqrt{1 - e^2}\,\sin E\bigr), \qquad
  \dot q = \frac{\bigl(-\sin E,\; \sqrt{1 - e^2}\cos E\bigr)}{1 - e\cos E} .
\end{equation*}
The energy $H = \tfrac12|\dot q|^2 - 1/|q|$ and the angular momentum
$L = q_1\dot q_2 - q_2\dot q_1$ are constant. Near the pericentre the force is
$1/(1 - e)^2 \approx 6$ times stronger than on average, which is what makes an
eccentric orbit demanding for a fixed step.
\begin{center}
\small
% Written by examples/kepler/main.f90: do not edit.
\begin{tabular}{ll}
\toprule
eccentricity & $e = 0.6$ (semi-major axis 1, period $2\pi$, $GM = 1$) \\
initial condition & $q = (0.4, 0)$, $\dot q = (0, 2)$: the pericentre \\
invariants & $H_0 = -0.5$, $L_0 = \sqrt{1 - e^2} = 0.8$ \\
solution figures & 200 steps per period, 10 periods \\
convergence & error $\max_i |u_i - u_i^{\mathrm{exact}}|$ after 3 periods, 100, 200, \dots, 12800 steps per period \\
observed order & slope of the two finest levels with errors in $(10^{-10}, 1)$ \\
long time & $\max|H - H_0|$, $\max|L - L_0|$ over 100 periods, 400 steps per period \\
exact solution & Kepler's equation $E - e\sin E = t$, Newton's method \\
methods & \texttt{rk3ws}, \texttt{rk4}, \texttt{dp54} \\
interface & plain arrays, \texttt{rk\_array\_scheme\_t}, state $(q_1, q_2, \dot q_1, \dot q_2)$ \\
\bottomrule
\end{tabular}

\end{center}

\subsection*{2. Solutions}
\begin{center}
\begin{tikzpicture}
  \begin{groupplot}[
      group style={group size=2 by 1, horizontal sep=1.9cm}, % two panels side by side
      width=6.6cm, height=6.2cm,                           % size of each panel
      grid=major, grid style={black!12},                   % light grid
    ]
    % ---- (a) orbits ----
    \nextgroupplot[
      title={(a) 10 orbits, 200 steps per period},         % panel title
      xlabel={$q_1$}, ylabel={$q_2$},                      % orbital plane
      axis equal image, xmin=-1.9, xmax=0.7, ymin=-1.1, ymax=1.1, % the ellipse
      legend to name=keplegend, legend columns=5,          % legend below both panels
      legend style={font=\small, /tikz/every even column/.append style={column sep=0.3cm}},
    ]
    \addplot[blue!80!black, thin] table[x=x, y=y] {data/orbit_rk4.dat};     % RK4: on the ellipse
    \addlegendentry{\texttt{rk4}}
    \addplot[black, thick, dashed] table[x=x, y=y] {data/ellipse.dat};      % exact ellipse
    \addlegendentry{exact}
    \addplot[only marks, mark=*, mark size=1.5pt] coordinates {(0, 0)};     % the central mass
    % ---- (b) position error ----
    \nextgroupplot[
      title={(b) position error, 200 steps per period},    % panel title
      xlabel={$t / 2\pi$ (periods)}, ylabel={$|q - q^{\mathrm{exact}}|$}, % error against periods
      ymode=log, ymin=1e-10, ymax=10, xmin=0, xmax=10,     % from dp54 to order one
    ]
    \addplot[green!50!black, dashed]  table[x expr=\thisrow{t}/(2*pi), y=error] {data/history_rk3ws.dat};    % order 2 here
    \addplot[blue!80!black]           table[x expr=\thisrow{t}/(2*pi), y=error] {data/history_rk4.dat};      % order 4
    \addplot[red!75!black]            table[x expr=\thisrow{t}/(2*pi), y=error] {data/history_dp54.dat};     % order 5
  \end{groupplot}
  \node[below=0.2cm] at (current bounding box.south)        % legend centred under both panels
    {\pgfplotslegendfromname{keplegend}};
\end{tikzpicture}
\end{center}
(a) At 200 steps per period the ellipse of \texttt{rk3ws} turns (precesses)
slowly; \texttt{rk4} stays on the exact ellipse at this scale. (b) The error grows mostly in jumps, one at each
pericentre passage (once per period), where the force changes fastest;
\texttt{rk3ws} (dashed) is the highest curve, \texttt{dp54} the lowest.

\newpage
\subsection*{3. Tests: convergence and the invariants $H$, $L$}

\paragraph{Prediction.} (i) The error after a fixed time converges at the
method's general order as the steps per period double. (ii) Explicit
Runge--Kutta methods keep neither the energy nor the angular momentum (a
quadratic invariant: only special implicit methods, such as the Gauss
methods, keep those), so both drift, by an amount that shrinks at the
method's order.

\paragraph{Result.}
\begin{center}
\small
% Written by examples/kepler/main.f90: do not edit.
\begin{tabular}{llcccc}
\toprule
method & name & order & observed & $\max|H - H_0|$ & $\max|L - L_0|$ \\
\midrule
RK3 of Wicker and Skamarock & \texttt{rk3ws} & 2 & 2.00 & $2.83\times10^{-3}$ & $1.89\times10^{-3}$ \\
Classic RK4 & \texttt{rk4} & 4 & 4.06 & $1.05\times10^{-5}$ & $1.85\times10^{-6}$ \\
Dormand-Prince 5(4) & \texttt{dp54} & 5 & 5.04 & $1.70\times10^{-7}$ & $1.66\times10^{-8}$ \\
\bottomrule
\end{tabular}

\end{center}

\begin{center}
\begin{tikzpicture}
  \begin{groupplot}[
      group style={group size=2 by 1, horizontal sep=1.9cm}, % two panels side by side
      width=6.6cm, height=6.2cm,                           % size of each panel
      grid=major, grid style={black!12},                   % light grid
      every axis plot/.append style={thick},               % same weight for all curves
    ]
    % ---- (i) convergence ----
    \nextgroupplot[
      xmode=log, ymode=log,                                % both axes logarithmic
      title={(i) error after 3 periods},                   % panel title
      xlabel={$\Delta t$},                                 % step
      legend to name=keptestlegend, legend columns=7,      % legend below both panels
      legend style={font=\small, /tikz/every even column/.append style={column sep=0.2cm}},
    ]
    \addplot[green!50!black, dashed, mark=triangle, mark size=1.5pt] table[x=dt, y=error] {data/error_rk3ws.dat}; % order 2 here
    \addlegendentry{\texttt{rk3ws}}
    \addplot[blue!80!black, mark=diamond*, mark size=1.5pt]   table[x=dt, y=error] {data/error_rk4.dat};      % order 4
    \addlegendentry{\texttt{rk4}}
    \addplot[red!75!black, mark=pentagon*, mark size=1.5pt]   table[x=dt, y=error] {data/error_dp54.dat};     % order 5
    \addlegendentry{\texttt{dp54}}
    % ---- (ii) energy drift ----
    \nextgroupplot[
      xmode=log, ymode=log,                                % both axes logarithmic
      title={(ii) $\max_{s \le t}|H - H_0|$, 400 steps/period}, % panel title: running maximum
      xlabel={$t / 2\pi$ (periods)}, ymin=1e-10,           % time in periods
    ]
    \addplot[green!50!black, dashed]  table[x expr=\thisrow{t}/(2*pi), y=energy] {data/invariants_rk3ws.dat};    % order 2 here
    \addplot[blue!80!black]           table[x expr=\thisrow{t}/(2*pi), y=energy] {data/invariants_rk4.dat};      % order 4
    \addplot[red!75!black]            table[x expr=\thisrow{t}/(2*pi), y=energy] {data/invariants_dp54.dat};     % order 5
  \end{groupplot}
  \node[below=0.2cm] at (current bounding box.south)        % legend centred under both panels
    {\pgfplotslegendfromname{keptestlegend}};
\end{tikzpicture}
\end{center}
(i) Every method converges at its general order once the step is small
enough (\texttt{rk3ws}: 2). The cut-off for the observed order is $10^{-10}$ here:
over tens of thousands of steps round-off accumulates above $10^{-12}$.
(ii) The energy drifts in steps, one per pericentre passage, and keeps
growing: no explicit Runge--Kutta method stays on the orbit.
\end{document}
